\documentclass[11pt]{article}
\usepackage[letterpaper,margin=1in]{geometry}
\usepackage{amsmath,amssymb,amsthm,mathtools}
\usepackage{graphicx}
\usepackage{xcolor}
\usepackage{enumitem}
\usepackage[colorlinks=true,linkcolor=blue!55!black,citecolor=blue!55!black,urlcolor=blue!55!black]{hyperref}

\newtheorem{theorem}{Theorem}
\newtheorem{lemma}[theorem]{Lemma}
\newtheorem{corollary}[theorem]{Corollary}
\newtheorem{proposition}[theorem]{Proposition}
\theoremstyle{definition}
\newtheorem{definition}[theorem]{Definition}
\theoremstyle{remark}
\newtheorem{remark}[theorem]{Remark}

\newcommand{\R}{\mathbb R}
\newcommand{\E}{\mathcal E_n}
\newcommand{\K}{\mathcal K}
\newcommand{\Span}{\operatorname{span}}
\newcommand{\Ext}{\operatorname{Ext}}
\newcommand{\diag}{\operatorname{diag}}
\newcommand{\sgn}{\operatorname{sgn}}
\newcommand{\one}{\mathbf 1}
\newcommand{\perpB}{\perp_B}

\title{Minimal Koldobsky Sets in $\ell_\infty^n$:\\
Exact Face Factorization, Nonvertex Frames, and the Two-Dimensional Classification}
\author{Automated mathematics run \texttt{fa\_banach\_001}, lane 12}
\date{August 9, 2026}

\begin{document}
\maketitle

\begin{center}
\fbox{\parbox{0.92\textwidth}{
\textbf{Status.} Strong partial result, likely valid.  We give an exact local
criterion at every cube-face point, classify all finite minimal $\K$-sets in
$\ell_\infty^2$, construct a continuum of nonvertex minimal $n$-point sets
for every $n\geq3$, classify all vertex-supported sets, and prove
$\kappa(\ell_\infty^n)=n$.  An intrinsic classification of all nonvertex
minimal sets for $n\geq3$ remains open.
}}
\end{center}

\section{Source question and results}

Sain, Manna, and Paul introduce $\K$-sets through local preservation of
Birkhoff--James orthogonality \cite{SMP}.  At the end of page 17 they ask:
\begin{quote}
\textbf{Open Question.} Characterize the minimal $\K$-sets in
$\ell_\infty^n$ and hence find $\kappa(\ell_\infty^n)$.
\end{quote}

\begin{figure}[ht]
\centering
\includegraphics[width=0.98\textwidth]{figures/open_problem_crop.png}
\caption{The source question on page 17 of arXiv:2407.08900.}
\end{figure}

For a real normed space $X$, write
\[
 x\perpB y\quad\Longleftrightarrow\quad
 \|x+\lambda y\|\geq\|x\|\quad(\lambda\in\R).
\]
A linear map $T:X\to X$ \emph{preserves orthogonality at $x$} if
$x\perpB y$ implies $Tx\perpB Ty$ for every $y$.  A subset $A\subset S_X$
is a $\K$-set if every linear map preserving at all $x\in A$ is a scalar
multiple of an isometry.  It is minimal if no proper subset is a $\K$-set.
The $\K$-number is the least cardinality of a finite $\K$-set.

We work in $X=\ell_\infty^n$.  The packet proves four main theorems.

\begin{theorem}[Arbitrary-face factorization]\label{thm:factor}
For $x\in S_{\ell_\infty^n}$ and $T\in\mathbb L(\R^n)$, suppose
$w=Tx\neq0$ and put $c=\|w\|_\infty$.  Then $T$ preserves orthogonality at
$x$ if and only if
\begin{equation}\label{eq:simplex-inclusion}
 J(x)\subset c^{-1}T^*J(w),
\end{equation}
where $J(u)=\{f\in S_{(\ell_\infty^n)^*}:f(u)=\|u\|_\infty\}$.
Equivalently, the extreme norming functionals of $x$ admit an entrywise
nonnegative row-stochastic factorization through the normalized pullbacks of
the extreme norming functionals of $w$.  If $Tx=0$, preservation at $x$ is
automatic.
\end{theorem}

\begin{theorem}[Near-vertex nonvertex frames]\label{thm:near}
Let $n\geq3$ and choose arbitrary $t_1,\ldots,t_n\in(-1,1)$.  Define
\[
 x^i_j=\begin{cases}t_i,&j=i,\\1,&j\neq i.\end{cases}
\]
Then $\{x^1,\ldots,x^n\}$ is a minimal $\K$-set in $\ell_\infty^n$.
In particular, every $n\geq3$ has a continuum of genuinely nonvertex
minimal $\K$-sets of optimal cardinality.
\end{theorem}

\begin{theorem}[Complete finite classification in dimension two]
\label{thm:two}
Up to replacing points by antipodes, the finite minimal $\K$-sets in
$\ell_\infty^2$ are exactly
\begin{enumerate}[label=\textup{(\roman*)}]
\item two linearly independent cube vertices;
\item one cube vertex, one smooth boundary point whose first coordinate is
active, and one smooth boundary point whose second coordinate is active.
\end{enumerate}
\end{theorem}

\begin{theorem}[Vertex sets and the exact $\K$-number]\label{thm:vertex}
Let $\E=\{-1,1\}^n$ and $A\subset\E$.  Then
\[
 A\text{ is a }\K\text{-set}\quad\Longleftrightarrow\quad
 \Span A=\R^n.
\]
Thus the minimal vertex-supported sets are exactly vertex bases, and for
every $n\geq2$,
\[
 \boxed{\ \kappa(\ell_\infty^n)=n\ }.
\]
\end{theorem}

The restriction ``finite'' in Theorem~\ref{thm:two} is intentional.  No
claim about possible infinite minimal sets is made.

\section{Idea of the proof}

At an arbitrary boundary point, Birkhoff--James orthogonality is a
straddling-zero condition on the finite list of active signed coordinate
functionals.  If the normalized pullback list at the image is strictly
positive on a direction, preservation forces the source list to have one
strict sign.  A segment from the distinguished point rules out the negative
sign.  Passing to the closed cone and applying Farkas duality produces the
exact stochastic factorization of norming simplices.

At a cube vertex the source simplex spans the whole dual space, so the local
factorization forces an invertible operator, a vertex image, and a
sign-conjugated nonnegative inverse.  With a vertex normalized to $\one$, the
inverse becomes row-stochastic.  Every additional test point then says
exactly that it lies in the stochastic image of the cube.

The nonvertex construction takes $n$ points just inside the $n$ different
coordinate directions from one vertex.  Their norming simplices overlap in
$n-2$ coordinates.  These overlaps first exclude singular preservers and
then equalize all image norms and all inverse-row $\ell_1$ norms.  Equality
in H\"older's inequality fixes signs on row supports.  Any shared support
coordinate becomes constant across all image columns, while each inverse row
needs a private coordinate because the corresponding diagonal parameter is
strictly below one.  There are exactly enough coordinates for one private
coordinate per row, forcing a monomial inverse.

In dimension two, a normalized vertex makes the inverse stochastic, and one
smooth test of each active-coordinate type makes one row at a time
deterministic.  This proves the three-point set; diagonal and triangular
stochastic deformations show that every missing finite type fails.  For
vertex bases in all dimensions, the positive inverse laws match source and
image row-sign patterns and again force a monomial operator.

\section{The exact arbitrary-face certificate}

For $u\in S_{\ell_\infty^n}$ let
\[
 I(u)=\{i:|u_i|=1\}.
\]
Let $F_u$ be the $|I(u)|\times n$ matrix whose row indexed by $i\in I(u)$
is $\sgn(u_i)e_i^*$.  The norming set is the simplex
\begin{equation}\label{eq:norming-simplex}
 J(u)=\operatorname{co}\{\sgn(u_i)e_i^*:i\in I(u)\}.
\end{equation}
James's criterion \cite{James} gives
\begin{equation}\label{eq:straddle-general}
 u\perpB z\quad\Longleftrightarrow\quad
 \min(F_uz)\leq0\leq\max(F_uz).
\end{equation}

\begin{proof}[Proof of Theorem~\ref{thm:factor}]
Only the case $w=Tx\neq0$ needs proof.  Set
\[
 A=F_x,\qquad B=c^{-1}F_wT.
\]
Every row of either matrix evaluates to one at $x$, so
$Ax=Bx=\one$.  By \eqref{eq:straddle-general}, preservation says that if
the coordinates of $Az$ straddle zero, then those of $Bz$ straddle zero.

We first claim
\begin{equation}\label{eq:strict-implication}
 Bz>0\quad\Longrightarrow\quad Az>0.
\end{equation}
Indeed, the contrapositive of preservation says that $Az$ is either strictly
positive or strictly negative.  In the latter case consider
$z(t)=(1-t)x+tz$.  The vector $Bz(t)$ remains strictly positive for
$0\leq t\leq1$, whereas $Az(t)$ moves from $\one$ to a strictly negative
vector and hence meets the straddling set.  This contradicts preservation
and proves \eqref{eq:strict-implication}.  Applying it to
$z+\varepsilon x$ and letting $\varepsilon\downarrow0$ yields
\begin{equation}\label{eq:closed-implication}
 Bz\geq0\quad\Longrightarrow\quad Az\geq0.
\end{equation}

For each row $a$ of $A$, implication \eqref{eq:closed-implication} says that
$a$ is nonnegative on the polyhedral cone $\{z:Bz\geq0\}$.  The Farkas
lemma identifies the dual cone with the nonnegative cone generated by the
rows of $B$.  Hence $A=CB$ for an entrywise nonnegative matrix $C$.
Evaluating at $x$ gives $\one=C\one$, so $C$ is row-stochastic.  Using
\eqref{eq:norming-simplex}, this is exactly
\eqref{eq:simplex-inclusion}.

Conversely, if $A=CB$ for a nonnegative row-stochastic $C$, then $Bz>0$
implies $Az>0$, and $Bz<0$ implies $Az<0$.  The contrapositive is precisely
the required straddling implication.  Finally, if $Tx=0$, the conclusion
$Tx\perpB Ty$ holds for every $y$ by definition.
\end{proof}

\begin{remark}
Theorem~\ref{thm:factor} turns local preservation into a finite linear
feasibility certificate.  It is exact, rather than merely a comparison of
the dimensions of the source and image faces.
\end{remark}

\section{Vertices and stochastic parallelotopes}

For a cube vertex $x$, put $D_x=\diag(x_1,\ldots,x_n)$.

\begin{corollary}[Local positive-inverse law]\label{cor:vertex-law}
Let $x\in\E$ and $Tx\neq0$.  Then $T$ preserves at $x$ if and only if there
are $c>0$ and $v\in\E$ such that
\[
 Tx=cv,\qquad T\text{ is invertible},\qquad
 cD_xT^{-1}D_v\geq0
\]
entrywise.
\end{corollary}

\begin{proof}
The simplex $J(x)$ contains $n$ linearly spanning extreme points.  Inclusion
\eqref{eq:simplex-inclusion} forces $T^*$ to be surjective and the image face
to have at least $n$ active coordinates.  Thus $T$ is invertible and every
coordinate of $Tx$ is active, so $Tx=cv$ for a vertex $v$.  In the matrix
factorization from Theorem~\ref{thm:factor},
\[
 D_x=C(c^{-1}D_vT),
\]
whence $C=cD_xT^{-1}D_v$.  Necessity and sufficiency now follow from the
factorization theorem.
\end{proof}

The vertex law has a useful global reformulation.  Suppose a test set
contains a vertex.  Source and output cube isometries and scalar
normalization reduce the nonzero case to
\begin{equation}\label{eq:stochastic-normalization}
 T\one=\one,\qquad R=T^{-1}\geq0,\qquad R\one=\one.
\end{equation}
Thus $R$ is invertible and row-stochastic.

\begin{proposition}[Stochastic-parallelotope criterion]\label{prop:stochastic}
Under \eqref{eq:stochastic-normalization}, for every
$x\in S_{\ell_\infty^n}$,
\[
 T\text{ preserves at }x
 \quad\Longleftrightarrow\quad
 \|R^{-1}x\|_\infty=1
 \quad\Longleftrightarrow\quad
 x\in R(B_{\ell_\infty^n}).
\]
\end{proposition}

\begin{proof}
Put $y=R^{-1}x=Tx$.  Since $R$ is row-stochastic,
$\|Ry\|_\infty\leq\|y\|_\infty$; hence $\|x\|_\infty=1$ shows that
$\|y\|_\infty\leq1$ is equivalent to $\|y\|_\infty=1$.

Assume $\|y\|_\infty=1$.  If $i\in I(x)$, equality in
\[
 x_i=\sum_jR_{ij}y_j
\]
forces $y_j=\sgn(x_i)$ wherever $R_{ij}>0$.  Consequently
$g_i=\sgn(x_i)\sum_jR_{ij}e_j^*$ belongs to $J(y)$, and
$T^*g_i=\sgn(x_i)e_i^*$.  Theorem~\ref{thm:factor} gives preservation.
Conversely, the same theorem represents every extreme functional
$\sgn(x_i)e_i^*$ as $c^{-1}T^*g_i$ with $g_i\in J(y)$ and
$c=\|y\|_\infty$.  Thus $g_i=cR^T\sgn(x_i)e_i^*$; its $\ell_1$ norm is
$c$ because row $i$ of $R$ has sum one.  Since $g_i$ is norming, its norm is
one, so $c=1$.
\end{proof}

Thus every vertex-containing instance reduces exactly to asking which test
sets force a stochastic image parallelotope $R(B_\infty)\subset B_\infty$
that contains them to arise from a permutation matrix.

\section{A continuum of nonvertex minimal sets}

\begin{proof}[Proof of Theorem~\ref{thm:near}]
Let $X=(x^1\ \cdots\ x^n)$ and
$D=\diag(1-t_1,\ldots,1-t_n)$.  Then $X=\one\one^T-D$.  The matrix
determinant lemma gives
\begin{equation}\label{eq:near-det}
 \det X=(-1)^n\prod_i(1-t_i)
 \left(1-\sum_i\frac1{1-t_i}\right).
\end{equation}
Every summand is greater than $1/2$, so the parenthesis is nonzero for
$n\geq3$.  Hence the $x^i$ form a basis.

Let $T$ preserve at every $x^i$.  Suppose first that $T$ is singular but
$Tx^i\neq0$ for some $i$.  The simplex inclusion in
Theorem~\ref{thm:factor} puts
\[
 H_i=\Span\{e_k^*:k\neq i\}
\]
inside $\operatorname{range}T^*$.  Therefore $\operatorname{rank}T=n-1$
and the range equals $H_i$.  If $Tx^j\neq0$ for a second index $j$, the same
argument gives $H_i=H_j$, impossible.  Thus at most one basis image is
nonzero and $\operatorname{rank}T\leq1$, contradicting $n-1\geq2$.  If all
images vanish then $T=0$.  We may therefore assume that a nonzero $T$ is
invertible.

Write $R=T^{-1}$, let $r_k$ be its $k$th row, and set
\[
 \rho_k=\|r_k\|_1,\qquad c_i=\|Tx^i\|_\infty.
\]
For $k\neq i$, $e_k^*$ is an extreme norming functional of $x^i$.
Invertibility and Theorem~\ref{thm:factor} say that its unique pulled-back
norming functional is $c_iR^Te_k^*$.  Taking its dual norm gives
\begin{equation}\label{eq:overlap-norms}
 c_i\rho_k=1\qquad(k\neq i).
\end{equation}
For any two indices there is a third; the overlaps in
\eqref{eq:overlap-norms} show that all $c_i$ equal a common $c$ and all
$\rho_k$ equal $c^{-1}$.  Rescale $T$ so that both common values are one.

Put $W=TX=(w^1\ \cdots\ w^n)$.  Then
\[
 RW=X,\qquad \|w^i\|_\infty=1,\qquad \|r_k\|_1=1.
\]
For $k\neq i$, the identity $r_k\cdot w^i=1$ is equality in H\"older's
inequality.  Hence
\begin{equation}\label{eq:holder-sign}
 r_{kj}\neq0, k\neq i
 \quad\Longrightarrow\quad w^i_j=\sgn(r_{kj}).
\end{equation}

Suppose output coordinate $j$ belongs to the supports of two different rows
$r_k,r_\ell$.  Choose $i$ different from both.  Equation
\eqref{eq:holder-sign} first shows that the two signs agree.  Applying the
same equation with row $k$ to every column other than $k$, and with row
$\ell$ to column $k$, shows that the entire $j$th row of $W$ is constant
with that sign.

Each row $r_k$ must now have a \emph{private} support coordinate, one which
belongs to no other row.  Otherwise every coordinate in its support is
shared, so the preceding observation gives
\[
 (RW)_{kk}=r_k\cdot w^k=\|r_k\|_1=1,
\]
contrary to $(RW)_{kk}=X_{kk}=t_k<1$.  Private coordinates for distinct rows
are distinct.  There are $n$ rows and only $n$ coordinates, so each row has
exactly one support coordinate and no coordinate is shared.  Thus $R$ is
monomial.  Since all row $\ell_1$ norms equal one, $R$ is a signed
permutation.  Undoing the normalization, $T$ is a scalar multiple of an
isometry.

Finally, every $\K$-set must span: if it does not, a nonzero rank-one map
annihilating its span preserves there and is not a scalar isometry.  Since
our set is an $n$-point basis, removal of any point destroys the spanning
property.  It is minimal.
\end{proof}

\section{Complete finite classification in the square}

Replacing $x$ by $-x$ does not change local preservation.  Projectively, the
square boundary consists of vertices, smooth points active in coordinate
one, and smooth points active in coordinate two.

\begin{proof}[Proof of Theorem~\ref{thm:two}]
First consider a triple of the stated form.  Normalize its vertex to
$\one=(1,1)$.  If $T\one=0$, write the columns of $T$ as $a,-a$.  A smooth
active-first point may be written $(1,u)$ with $|u|<1$ and satisfies
$(1,u)\perpB e_2$.  Preservation requires $(1-u)a\perpB-a$, which forces
$a=0$.  Thus $T=0$.

If $T\one\neq0$, normalize as in \eqref{eq:stochastic-normalization}, so
$R=T^{-1}$ is invertible and row-stochastic.  At a smooth point active in
coordinate $i$, Proposition~\ref{prop:stochastic} gives $x=Ry$ for some
$y\in B_\infty$.  Equality in active row $i$ forces $y$ to have the active
sign on the support of that row.  If both entries of the row were positive,
$y$ would be a constant vertex; every stochastic row would then give the
same constant vertex $x$, contradicting smoothness.  Hence row $i$ is a
coordinate row.  One smooth point of each active type forces both rows of
$R$ to be distinct coordinate rows.  Therefore $R$ is a permutation and
$T$ is a scalar cube isometry.  The triple is a $\K$-set.

We next show that the list is exhaustive and minimal.  A finite set of
smooth points only is preserved by
\[
 T=\diag(1,r)
\]
for some $r>1$ sufficiently close to one: all active-first points retain
first-coordinate activity, and all active-second points retain
second-coordinate activity.  This $T$ is not a scalar isometry.

After normalizing a vertex to $\one$, a finite collection of smooth points
of active type one may be written projectively as $(1,u_s)$, $|u_s|<1$.
Choose
\[
 0<a<\min_s\frac{1+u_s}{2},\qquad
 R=\begin{pmatrix}1&0\\a&1-a\end{pmatrix}.
\]
Then
\[
 R^{-1}(1,u_s)^T=\left(1,\frac{u_s-a}{1-a}\right)^T\in B_\infty.
\]
Proposition~\ref{prop:stochastic} says that the nonisometry $R^{-1}$
preserves at the vertex and at all these points.  The transposed construction
handles active type two.

Finally, two nonantipodal projective vertex directions are linearly
independent and hence form a $\K$-set by Theorem~\ref{thm:vertex}, whose
proof appears below.  A finite minimal set containing two such directions is
therefore exactly that pair.  A minimal set with only one vertex direction
must contain both smooth types, while any one point of each type with the
vertex already forms the triple proved sufficient above.  Antipodal repeats
are redundant.  These observations prove both exhaustiveness and minimality.
\end{proof}

\section{Vertex-supported classification and the \texorpdfstring{$\K$}{K}-number}

We include the vertex proof because it both supplies Theorem~\ref{thm:two}
and determines the numerical invariant in the source question.

\begin{lemma}[Common image norm]\label{lem:common-norm}
If $x,y\in\E$ and $T$ preserves at both points, then
$\|Tx\|_\infty=\|Ty\|_\infty$.
\end{lemma}

\begin{proof}
The conclusion is immediate for $y=\pm x$.  Otherwise, $x$ and $y$ agree in
some coordinate and disagree in another, so the vertex straddling criterion
gives $x\perpB(x+y)$ and $y\perpB(x+y)$.  If one image vanishes,
preservation forces the other to be orthogonal to itself and hence to vanish.
Otherwise Corollary~\ref{cor:vertex-law} writes
$Tx=cv$ and $Ty=dw$ for $c,d>0$ and vertices $v,w$.  Preservation gives
\[
 cv\perpB(cv+dw),\qquad dw\perpB(cv+dw).
\]
Choosing scalar $-1$ in the definition of orthogonality in the first
relation yields $d\geq c$; the second yields $c\geq d$.
\end{proof}

\begin{proof}[Proof of Theorem~\ref{thm:vertex}]
If $A$ does not span, choose a nonzero functional $f$ annihilating its span
and a nonzero vector $u$.  The rank-one map $Tz=f(z)u$ vanishes on $A$, hence
preserves there, but is not a scalar isometry for $n\geq2$.

Conversely, choose a vertex basis $x^1,\ldots,x^n$ from a spanning $A$ and
suppose $T$ preserves there.  By Lemma~\ref{lem:common-norm}, their image
norms have a common value $c$.  If $c=0$, then $T=0$.  Otherwise
Corollary~\ref{cor:vertex-law} gives vertices $v^i$ with
$Tx^i=cv^i$, makes $T$ invertible, and gives
\begin{equation}\label{eq:many-positive-new}
 cD_{x^i}T^{-1}D_{v^i}\geq0\qquad(i=1,\ldots,n).
\end{equation}
Let $X=(x^1\ \cdots\ x^n)$, $V=(v^1\ \cdots\ v^n)$, and
$R=T^{-1}=(r_{ab})$.  Both $X$ and $V=c^{-1}TX$ are invertible.  If
$r_{ab}\neq0$, inequalities \eqref{eq:many-positive-new} imply
\begin{equation}\label{eq:row-match-new}
 (v_b^1,\ldots,v_b^n)
 =\sgn(r_{ab})(x_a^1,\ldots,x_a^n).
\end{equation}
Two nonzero entries in column $b$ of $R$ would make two rows of $X$
proportional, contrary to invertibility.  Every column therefore contains
exactly one nonzero entry, and invertibility puts them in distinct rows.
Thus $R$ and $T$ are monomial.  Since $Tx^i=cv^i$, every nonzero monomial
entry of $T$ has absolute value $c$.  Hence $T$ is a scalar signed
permutation.  This proves the equivalence, and inclusion-minimal spanning
vertex sets are precisely bases.

The same rank-one annihilator shows that every finite $\K$-set, not only a
vertex set, has at least $n$ points.  Vertex bases exist; for example, the
columns defined by $x^i_j=-1$ for $j<i$ and $x^i_j=1$ for $j\geq i$ form a
basis.  The lower bound is attained, proving
$\kappa(\ell_\infty^n)=n$.
\end{proof}

\section{A sharp obstruction to the naive coordinate test}

One might hope that $\{\one,e_1,\ldots,e_n\}$ is a $\K$-set.  Already for
$n=3$ this fails.  Let $J$ be the all-ones matrix and put
\begin{equation}\label{eq:obstruction}
 T=J-2I,\qquad T^{-1}=\tfrac12(J-I).
\end{equation}
The inverse is nonnegative and row-stochastic, and $T\one=\one$, so
Corollary~\ref{cor:vertex-law} gives preservation at $\one$.  Moreover,
$Te_i$ has $-1$ in coordinate $i$ and $1$ in the other two coordinates, and
\[
 g_i=\tfrac12\sum_{j\neq i}e_j^*\in J(Te_i),
 \qquad T^*g_i=e_i^*.
\]
Theorem~\ref{thm:factor} gives preservation at every $e_i$.  But $T$ is not
a scalar signed permutation.  Taking the block sum of this $3\times3$
matrix with an identity gives the same obstruction in all $n>3$.

This counterexample explains why Theorem~\ref{thm:near} uses points with
$n-1$ overlapping active coordinates: one smooth coordinate test in each
direction does not give enough rigidity.

\section{Verification, literature boundary, and limitations}

The script
\begin{center}
\texttt{code/verify\_cube\_vertex\_k\_sets.py}
\end{center}
implements Theorem~\ref{thm:factor} as a linear feasibility test.  It
retains the exhaustive vertex-image enumeration through dimension four,
the 8,800-point positive-inverse grid check, and 1,000 random vertex-basis
checks.  It also exhausts all $3^9=19{,}683$ ternary matrices on the
near-vertex frame
\[
 (0,1,1),\quad(1,0,1),\quad(1,1,0),
\]
finding precisely the $48$ nonzero signed permutations; checks 200 random
near-vertex frames through dimension seven; exhausts 625 small integer
matrices on a representative two-dimensional triple and checks
nonisometries for its proper pair types; and verifies
\eqref{eq:obstruction}.  It terminates with
\begin{center}
\texttt{ALL UNRESTRICTED K-SET PROGRAM CHECKS PASSED}.
\end{center}
These are stress tests, not substitutes for the analytic proofs.

Bounded searches on August 9, 2026 used the run indexes and arXiv with
``minimal K-sets ell-infinity'', ``Koldobsky-set'', ``K-number orthogonality
isometries'', and close variants.  The source \cite{SMP} and the later broad
refinement \cite{MMPS} were inspected.  No exact statement of the
arbitrary-face factorization, near-vertex theorem, dimension-two
classification, vertex-spanning theorem, or exact $\K$-number was found.
This is a bounded novelty assessment, not an exhaustive review.

The factorization theorem is a complete local certificate and reduces every
finite candidate to explicit nonnegative matrix constraints.  It has not yet
been converted into a closed intrinsic classification of all finite
nonvertex minimal sets for $n\geq3$.  Possible infinite minimal sets are also
outside the claim.  These are the precise remaining boundaries.

Human review should focus on the Farkas step in
Theorem~\ref{thm:factor}, the singular case and private-coordinate argument
in Theorem~\ref{thm:near}, and the finite-set qualifier in
Theorem~\ref{thm:two}.  Recommended status: \emph{strong partial result,
likely valid}.

\begin{thebibliography}{9}

\bibitem{SMP}
D. Sain, J. Manna, and K. Paul,
\emph{On local preservation of orthogonality and its application to
isometries}, Linear Algebra Appl. 690 (2024), 112--131;
arXiv:2407.08900.

\bibitem{James}
R. C. James,
\emph{Orthogonality and linear functionals in normed linear spaces},
Trans. Amer. Math. Soc. 61 (1947), 265--292.

\bibitem{MMPS}
J. Manna, K. Mandal, K. Paul, and D. Sain,
\emph{Refinements of the Blanco--Koldobsky--Turn\v{s}ek Theorem},
Electron. J. Linear Algebra 41 (2025), 616--631;
arXiv:2504.11161.

\end{thebibliography}

\end{document}
